\begin{graphicspathcontext}{{./chapters/ai/imgs/},{./chapters/ai/imgs/auto/},\old}

\begin{frame}[c]{Qu'est-ce que l'IA ?}
	\begin{definition}[axée sur le mimétisme comportemental, par Minsky, années 50]
		L'IA est la science de la programmation des ordinateurs pour \emph{accomplir des tâches} qui nécessitent de \emph{l'intelligence lorsqu'elles sont réalisées par des êtres humains}
	\end{definition}
	\vspace{1cm}
	\begin{definition}[de l'OCDE]
		Une IA est un système fondé sur des machines qui, pour des \emph{objectifs explicites ou implicites}, déduit à partir des \emph{entrées} qu'il reçoit comment générer des résultats tels que des \emph{prévisions, du contenu, des recommandations ou des décisions} susceptibles d'influencer les environnements physiques ou virtuels. Les \emph{différents systèmes d'IA} varient dans leurs niveaux d'autonomie et d'adaptabilité une fois déployés
	\end{definition}
\end{frame}

\sidecite{ElFallahSeghrouchni18}
\begin{frame}[c]{Vision holistique de l'intelligence (humaine ou artificielle)}
	\begin{columns}
		\begin{column}{.5\linewidth}
			\begin{definitionblock}{Intelligence \cite{piaget1947psychologie}}
				L'intelligence, c'est ce qu'on utilise quand on ne sait pas quoi faire : quand ni l'inné ni l'apprentissage ne nous a préparé à la situation particulière
			\end{definitionblock}
		\end{column}
		\begin{column}{.5\linewidth}
			\viconbox*[.49\linewidth]{Intelligence cognitive}{intelligences-cognitive}
			\hfill
			\viconbox*[.49\linewidth]{Intelligence collective}{intelligences-collective} \\[.5cm]
			\viconbox*[.49\linewidth]{Intelligence\\sociale}{intelligences-social}
			\hfill
			\viconbox*[.49\linewidth]{Intelligence émotionnelle}{intelligences-emotional}
		\end{column}
	\end{columns}
\end{frame}

\sidecite{ElFallahSeghrouchni18}
\figureslide{{Différentes approches} pour l'IA}{ai_approaches}

\sidecite{ElFallahSeghrouchni18}
\animatedfigureslide{{Les approches de l'IA} plus en détail}{ai_approaches_in_details_fr}

\end{graphicspathcontext}

\endinput

